\lyseis{Β1.3}{Μέτρο γωνίας}
\sol{1} Ελεύθερη κατασκευή. Η $\hat{\gamma}$ είναι ορθή και η $\hat{\varepsilon}$ έχει πλευρές αντικείμενες ημιευθείες. Έλεγξε κάθε γωνία μετρώντας ξανά.
\sol{2} Ελεύθερη κατασκευή. Και οι δύο έχουν μέτρο $40^\circ$, άρα είναι ίσες: το μέτρο εξαρτάται μόνο από το άνοιγμα και όχι από το μήκος των πλευρών.
\sol{3} α) $12{,}5^\circ$, γιατί $30'=0{,}5^\circ$. β) $40{,}25^\circ$, γιατί $15'=0{,}25^\circ$. γ) $18^\circ45'$, γιατί $0{,}75\cdot60=45$. δ) $1^\circ=60'=3600''$.
\sol{4} Σε μοίρες: $45$, $40{,}5$, $40{,}25$, $45{,}17$. Σειρά: $40{,}25^\circ<40^\circ30'<45^\circ<45^\circ10'$.
\sol{5} α) $x\hat{\mathrm{Ο}}z=35^\circ+48^\circ=83^\circ$. β) $y\hat{\mathrm{Ο}}z=120^\circ-75^\circ=45^\circ$.
\sol{6} Ελεύθερη κατασκευή. Σε κάθε τρίγωνο το άθροισμα των τριών γωνιών βγαίνει περίπου $180^\circ$ (με μικρές αποκλίσεις λόγω μέτρησης).
\sol{7} α) Σωστό. β) Λάθος: το μέτρο δεν εξαρτάται από το μήκος των πλευρών. γ) Σωστό. δ) Λάθος: $1^\circ=60'$. ε) Σωστό. στ) Σωστό.
\sol{8} Συνολική στροφή $45^\circ+30^\circ=75^\circ$. Μετά την αντίθετη στροφή απομένει $75^\circ-20^\circ=55^\circ$ από την αρχική πορεία.
\sol{9} Κάθε ώρα αντιστοιχεί σε $360^\circ:12=30^\circ$. Στις $2{:}00$: $60^\circ$. Στις $5{:}00$: $150^\circ$. Στις $6{:}00$: $180^\circ$.
\sol{10} $84^\circ-50^\circ=34^\circ$.
\sol{11} $35^\circ+50^\circ=85^\circ$ και μετά $85^\circ-20^\circ=65^\circ$ ανοιχτή.
\sol{12} Η γωνία φαίνεται μικρότερη από την ορθή, άρα το μέτρο της είναι μικρότερο από $90^\circ$: το σωστό είναι $50^\circ$. Οι δύο κλίμακες δίνουν άθροισμα $130^\circ+50^\circ=180^\circ$.
\sol{13} Οι πλευρές μιας γωνίας είναι ημιευθείες και δεν έχουν μήκος· το μέτρο δείχνει μόνο πόσο ανοίγουν. Αν μακρύνουμε ή κονταίνουμε τη γραμμή που σχεδιάσαμε, το άνοιγμα δεν αλλάζει.
\sol{14} $x\hat{\mathrm{Ο}}t=70^\circ+50^\circ+40^\circ=160^\circ$.
\sol{15} Αντιγράφουμε τη μία γωνία σε διαφανές χαρτί και την τοποθετούμε πάνω στην άλλη με κοινή κορυφή και μία κοινή πλευρά. Αν οι άλλες πλευρές συμπίπτουν, είναι ίσες. Αλλιώς, η γωνία που μένει στο εσωτερικό της άλλης είναι η μικρότερη.
\sol{16} Ελεύθερη κατασκευή. α) $x\hat{\mathrm{Ο}}t=60^\circ$. β) $x\hat{\mathrm{Ο}}t=15^\circ$. Έλεγξε με το μοιρογνωμόνιο.
\sol{17} Ελεύθερη απάντηση: π.χ. η ανοιχτή πόρτα $\approx$ μερικές δεκάδες μοίρες, η γωνία του τραπεζιού $90^\circ$.
\sol{18} Ελεύθερη κατασκευή. Όσο ανοίγει ο διαβήτης, τόσο μεγαλώνει το μέτρο της γωνίας των σκελών του.
